% SNN 信用分配：从原公式到在线/局部化近似
\documentclass[journal,12pt,onecolumn]{IEEEtran}
\IEEEoverridecommandlockouts
\usepackage{cite}
\usepackage{amsmath, amssymb, bm, graphicx, enumitem}
\usepackage{algorithmic}
\usepackage{algorithm}
\usepackage{booktabs}
\usepackage{textcomp}
\usepackage{xcolor}
\usepackage{tabularx}
\usepackage{makecell}
\usepackage{tikz}
\usetikzlibrary{arrows.meta, positioning}
\usepackage[UTF8]{ctex}

\input{math_commands_TESS.tex}
\input{math_commands_STLLR.tex}

\definecolor{myred}{RGB}{255, 67, 20}
\definecolor{myblue}{RGB}{51, 51, 255}
\definecolor{mygreen}{RGB}{0, 128, 0}

\begin{document}

\title{SNN 信用分配：从原公式到在线/局部化近似}

\author{}
\date{}

\maketitle

\begin{abstract}
本方案以六篇与 SNN 在线/局部信用分配密切相关的核心论文为理论主线，从每篇的\textbf{原公式}出发，逐步推导它们各自的梯度计算方式。

\textbf{研究基准}：从 BPTT 与 RTRL 两类精确信用分配基准出发（reverse-mode / forward-mode），分析在线/局部算法在 \textbf{temporal representation}、\textbf{eligibility} 与 \textbf{learning signal} 上进行了何种重组或近似。

\textbf{总研究问题}：What information is lost by temporal/spatial localization, and how do the losses interact?

三条研究路径：
\begin{itemize}[leftmargin=1.5em,itemsep=0pt]
\item \textbf{时间近似}：BPTT $\to$ OTTT $\to$ NDOT，研究 temporal dependency $\epsilon^l[t]$ 如何被 $\lambda I$ 或 $\operatorname{diag}(e^l[t])$ 替代；
\item \textbf{精确到局部}：RTRL $\to$ e-prop，研究 RTRL 的 full forward sensitivity 与 e-prop 的 local eligibility + learning-signal factorization 之间的关系；
\item \textbf{生物局部化}：STDP $\to$ S-TLLR $\to$ TESS，研究预突触-突触后时序信息如何逐步引入 task credit 与空间局部性。
\end{itemize}

核心研究问题：这些在线/局部化操作到底保留了什么 credit information，丢失了什么 credit information，以及这些差异什么时候真正影响学习。
\end{abstract}

\section*{文章汇总}

RTRL: \textbf{A learning algorithm for continually running fully recurrent neural networks} \cite{williams1989learning};

E-prop: \textbf{A solution to the learning dilemma for recurrent networks of spiking neurons} \cite{bellec_solution_2020};

OTTT: \textbf{Online Training Through Time for Spiking Neural Networks} \cite{xiao_online_2022};

NDOT: \textbf{NDOT: Neuronal Dynamics-based Online Training for Spiking Neural Networks} \cite{jiang_ndot_nodate};

S-TLLR: \textbf{S-TLLR: STDP-inspired Temporal Local Learning Rule for Spiking Neural Networks} \cite{apolinario_s-tllr_2024};

TESS: \textbf{TESS: A Scalable Temporally and Spatially Local Learning Rule for Spiking Neural Networks} \cite{apolinario_tess_2025}.

\section{研究背景与核心问题}

\subsection{SNN 训练的根本困难}

SNN 用二进制脉冲 $s_j^t\in\{0,1\}$ 传递信息：神经元只有在膜电位达到阈值时才发放一个脉冲，脉冲本身是二值的、不可微的。这带来两个层面的信用分配困难：

\begin{enumerate}[leftmargin=1.5em,itemsep=0pt]
\item \textbf{时间信用分配}：一个输出时刻的损失，可能由多个时间步之前的输入决定；要把这个损失归因到具体某个时刻的突触权重上，需要追踪"过去 $\to$ 现在"的因果链；
\item \textbf{空间信用分配}：损失发生在输出层，但要归因到每个隐藏层的权重上；需要把误差从输出层反向传播到各隐藏层。
\end{enumerate}

BPTT 是求解这两个问题的标准方法。它的核心思想是：把 SNN 沿时间展开成一个很深的 feedforward 计算图，然后在这个展开图上调用反向传播。代价是需要保存所有时间步的 forward states（每个时刻的膜电位 $u^l[t]$、脉冲 $s^l[t]$ 等），内存 $O(TLn)$，且必须等整个序列前向完成后才能开始反向，这排除了在线学习。

\textbf{关于"精确"的约定}：BPTT 与 RTRL 对一个\textbf{可微} recurrent computational graph 分别给出 reverse-mode 和 forward-mode exact gradient。但 SNN 的脉冲发放 $s=H(u-V_{th})$ 几乎处处导数为零，不是可微函数；实际训练中所有方法都用 surrogate $\Psi'_{\mathrm{SG}}$ 替代 $\partial s/\partial u$。因此本文所用的 reference 是 \textbf{BPTT-with-SG / RTRL under the same surrogate rule}，其"精确"仅指相对于选定 surrogate backward rule 的精确计算，不代表连续动力学的真实导数。下文所有 $g_{\mathrm{BPTT}}$ 均指 \textbf{BPTT-with-SG reference estimator}，不称为"exact gradient"。

\subsection{已有工作}

\textbf{时间近似主线}：OTTT \cite{xiao_online_2022} 注意到在 BPTT 的 temporal chain 里，spike-mediated path 通常极稀疏（因为 $H$ 几乎处处导数为零），于是丢弃该路径，把时序依赖退化为固定指数衰减 $\lambda^{t-\tau}$；NDOT \cite{jiang_ndot_nodate} 进一步从连续 LIF 微分方程推导动态比率 $e^l[t]$，替代固定 $\lambda$。

\textbf{精确到局部主线}：RTRL \cite{williams1989learning} 用 forward-mode 计算参数敏感度 $P_t$，是完全 forward-in-time 的精确梯度方法，但内存 $O(n^3)$；e-prop \cite{bellec_solution_2020} 重新组织 chain rule，把梯度分解为"局部 eligibility $\times$ learning signal"，大幅降低内存。

\textbf{生物局部化主线}：STDP 用 pre/post 时序定义突触可塑性；S-TLLR \cite{apolinario_s-tllr_2024} 引入 STDP 形式的因果与非因果资格迹，再用 learning signal 调制；TESS \cite{apolinario_tess_2025} 通过本地学习信号生成实现时空全局部。

\textbf{关于"联合局部化"的现状}：时空联合局部学习本身已经被研究过——TESS \cite{apolinario_tess_2025} 已把 temporal eligibility 与 layer-local spatial signal 结合；Traces Propagation 等工作也处理时空联合的 forward-only learning。因此本研究不以"首次实现时空联合"作为 novelty。

\textbf{本研究的问题不是"能否联合"，而是"联合局部化时各自损失了什么、损失如何耦合"}：
\begin{itemize}[leftmargin=1.5em,itemsep=0pt]
\item 时间局部化（例如 OTTT 的固定 $\lambda$、NDOT 的 $\operatorname{diag}(e^l)$）在 temporal dependency 上引入了何种结构性误差？
\item 空间局部化（例如 Mostafa-style local classifier）在目标替换上引入了何种结构性误差？
\item 当两种局部化同时施加时，误差是简单叠加还是产生 interaction？
\end{itemize}

\subsection{核心研究问题}

\textbf{总问题}：What information is lost by temporal/spatial localization, and how do the losses interact?

\textbf{三条研究路径下的具体问题}：
\begin{itemize}[leftmargin=1.5em,itemsep=0pt]
\item \textbf{路径一}：$\epsilon^l[t]\to\lambda I$ 或 $\operatorname{diag}(e^l[t])$ 到底改变了多少 temporal credit？在什么样的 $T,\lambda$ 下差异显现？
\item \textbf{路径二}：RTRL 的 full forward sensitivity 与 e-prop 的 local eligibility + learning signal factorization 是同维度对象吗？它们各自保留了什么？
\item \textbf{路径三}：STDP-inspired eligibility 中的 causal / non-causal timing 分别提供什么？layer-local spatial learning signal 与 cross-layer BP learning signal 的差别落在哪？
\end{itemize}

\section{数学基础}

\subsection{LIF 神经元（feedforward 形式）}

\begin{equation}
u_j^l[t+1]=\lambda\big(u_j^l[t]-V_{th}\,s_j^l[t]\big)+\sum_i W_{ji}^{l\leftarrow l-1}\,s_i^{l-1}[t+1]+b_j^l,
\qquad
s_j^l[t]=H(u_j^l[t]-V_{th}).
\end{equation}

\textbf{这个公式在说什么}：第 $l$ 层神经元 $j$ 的膜电位 $u_j^l[t+1]$ 由三部分决定——上一时刻膜电位经过衰减后剩下的"记忆"、来自上一层神经元的加权输入、以及偏置。当膜电位达到阈值 $V_{th}$ 时，神经元发放一个脉冲 $s_j^l[t]=1$。

\textbf{逐项拆解}：

\begin{itemize}[leftmargin=1.5em,itemsep=0pt]
\item $\lambda(u_j^l[t]-V_{th}s_j^l[t])$：这是\textbf{leak}加\textbf{reset}。
\begin{itemize}[leftmargin=1.2em,itemsep=0pt]
\item $u_j^l[t]$ 乘 $\lambda$ 表示膜电位随时间指数衰减（"leak"）。$\lambda=e^{-\delta t/\tau_m}<1$，$\delta t$ 越小或 $\tau_m$ 越大，$\lambda$ 越接近 1，膜电位衰减越慢、记忆越长；
\item $-V_{th}s_j^l[t]$ 是\textbf{reset by subtraction}：如果 $t$ 时刻发放了脉冲（$s_j^l[t]=1$），则从上一步膜电位中减去 $V_{th}$，防止膜电位一直累积。整个括号项先做 reset、再被 $\lambda$ 衰减（注意：$\lambda$ 作用在 reset 之后的量上）。
\end{itemize}
\item $\sum_i W_{ji}^{l\leftarrow l-1}s_i^{l-1}[t+1]$：来自上一层的加权输入。$s_i^{l-1}[t+1]$ 是上一层神经元 $i$ 在 $t+1$ 时刻的脉冲；如果是 feedforward 结构，输入到达的时刻和神经元更新同步。
\item $b_j^l$：偏置项，用来调整该神经元的兴奋程度基线。
\item $s_j^l[t]=H(u_j^l[t]-V_{th})$：Heaviside 阶跃函数 $\mathbb I[x\ge 0]$。当膜电位达到或超过阈值时输出 1，否则输出 0。这一项在训练中\textbf{前向不可导}，是 SNN 训练困难的核心来源。
\end{itemize}

\textbf{参数解释}：

\begin{center}
\begin{tabular}{lll}
\toprule
符号 & 数学定义 / 形状 & 物理含义与典型取值 \\
\midrule
$l\in\{1,\dots,L\}$ & 整数层索引 & $l=0$ 为输入层，$l=L$ 为输出层 \\
$j$ & 当前层神经元索引 & 遍历 $1,\dots,n^{(l)}$ \\
$i$ & 上一层神经元索引 & BPTT--OTTT--NDOT 主线中只指 $l-1$ 层 \\
$u_j^l[t]$ & $\mathbb R$（标量） & 第 $l$ 层 $j$ 号神经元 $t$ 时刻的膜电位，是 hidden state \\
$s_j^l[t]$ & $\{0,1\}$（标量） & 是否发放：$u_j^l[t]\ge V_{th}$ 时为 1 \\
$W_{ji}^{l\leftarrow l-1}$ & $\mathbb R^{n^{(l)}\times n^{(l-1)}}$ & 从 $l-1$ 层 $i$ 到 $l$ 层 $j$ 的突触权重 \\
$b_j^l$ & $\mathbb R$ & 偏置项 \\
$V_{th}$ & 正实数 & 发放阈值，本文取 $V_{th}=1$ \\
$\lambda$ & $e^{-\delta t/\tau_m}\in(0,1)$ & 膜衰减因子；$\tau_m$ 为膜时间常数，$\delta t$ 为离散步长 \\
$H(x)$ & $\mathbb I[x\ge 0]$ & Heaviside 阶跃函数，前向不可导 \\
\bottomrule
\end{tabular}
\end{center}

\subsection{Surrogate Gradient}

\begin{equation}
\frac{\partial s_j^l[t]}{\partial u_j^l[t]}\;\to\;\Psi'_{\mathrm{SG}}(u_j^l[t]-V_{th}),
\qquad
\Psi'_{\mathrm{SG}}(x)=\frac{1}{(1+\beta|x|)^2},\quad \beta=4.
\end{equation}

\textbf{为什么需要 surrogate}：Heaviside 函数 $H(x)$ 在 $x=0$ 处不可导，在其他地方导数为零。如果严格按导数反向传播，梯度将处处为零，训练无法进行。surrogate gradient 的做法是：\textbf{前向仍然使用真正的 Heaviside 产生脉冲，反向把 $\partial s/\partial u$ 替换成一个光滑、有界的函数}，让梯度能够回流。

\textbf{为什么选 fast-sigmoid}：$\Psi'_{\mathrm{SG}}(x)=1/(1+\beta|x|)^2$ 在 $x=0$ 处取最大值 1，随着 $|x|$ 增大迅速衰减。它有两个好处：
\begin{itemize}[leftmargin=1.5em,itemsep=0pt]
\item 当膜电位接近阈值（$u-V_{th}\approx 0$）时，神经元处于"决策边界"，梯度最大，学习信号最强；
\item 当膜电位远离阈值（$|u-V_{th}|$ 很大）时，梯度快速衰减到零，避免已经在"确定发放/确定不发放"状态的神经元继续被扰动。
\end{itemize}

\textbf{参数解释}：
\begin{itemize}[leftmargin=1.5em,itemsep=0pt]
\item $\Psi'_{\mathrm{SG}}$：surrogate derivative，用于反向替代不可导的 $\partial s/\partial u$；
\item $\beta$：锐度参数，本文取 $\beta=4$；$\beta$ 越大，surrogate 越接近真正的阶跃导数（几乎处处为零），训练信号越稀疏；
\item $D_{\Psi,t}=\operatorname{diag}[\Psi'_{\mathrm{SG}}(u_j^l[t]-V_{th})]_j$：把所有神经元 $j$ 的 surrogate 导数排成对角矩阵，$D_{\Psi,t}\in\mathbb R^{n^{(l)}\times n^{(l)}}$；
\item 该 surrogate 是本文的实验选择，不是原论文的统一规定。选择它只是为了在实验对照中固定一个明确的 surrogate 形式。
\end{itemize}

\subsection{BPTT：时间反向传播（原论文形式）}

\textbf{Step 1：ordered temporal product。}
\begin{equation}
\mathcal{E}_{t,\tau}
:=
\epsilon[t-1]\,\epsilon[t-2]\cdots\epsilon[\tau],
\qquad \tau<t.
\end{equation}

\textbf{解释}：从时间步 $\tau$ 的 state 到时间步 $t$ 的 state，需要经过一连串状态转移。链式法则要求把每一步的 Jacobian $\epsilon[i]=\partial u[i+1]/\partial u[i]$ 按时间顺序乘起来。注意矩阵乘法不可交换：从 $\tau$ 到 $t$ 的传播必须先应用 $\epsilon[\tau]$、再 $\epsilon[\tau+1]$、……、最后 $\epsilon[t-1]$，所以写成从 $t-1$ 到 $\tau$ 的下降顺序。$\mathcal{E}_{t,t}=I$（空乘积约定）。

\textbf{Step 2：BPTT 主公式。}

对权重 $W^{l\leftarrow l-1}$，损失 $L$ 的梯度由两部分组成：
\begin{equation}
\frac{\partial L}{\partial W^{l\leftarrow l-1}}
=
\sum_{t=1}^{T}
\frac{\partial L}{\partial s^{l}[t]}
\frac{\partial s^{l}[t]}{\partial u^{l}[t]}
\left(
\frac{\partial u^{l}[t]}{\partial W^{l\leftarrow l-1}}
+
\sum_{\tau<t}
\mathcal{E}_{t,\tau}
\frac{\partial u^{l}[\tau]}{\partial W^{l\leftarrow l-1}}
\right).
\end{equation}

\textbf{这个公式在说什么}：对每一时刻 $t$，权重 $W^{l\leftarrow l-1}$ 影响 $u^l[t]$ 有两条路径：
\begin{itemize}[leftmargin=1.5em,itemsep=0pt]
\item \textbf{直接影响}（第一项 $\partial u^l[t]/\partial W^{l\leftarrow l-1}$）：当前时刻 $t$，权重的输入直接进入 $u^l[t]$；
\item \textbf{历史路径}（第二项）：权重在某个过去时刻 $\tau$ 注入的输入，经过 $t-\tau$ 次状态转移传到现在。$\mathcal{E}_{t,\tau}$ 就是这个传播链。
\end{itemize}

外面乘的 $\partial L/\partial s^l[t]\cdot\partial s^l[t]/\partial u^l[t]$ 把 $u^l[t]$ 对 $L$ 的贡献转换成"该时刻该层神经元收到的 error signal"。

\textbf{temporal dependency 的定义}：
\begin{equation}
\epsilon^{l}[i]
=
\frac{\partial u^{l}[i+1]}{\partial u^{l}[i]}
+
\frac{\partial u^{l}[i+1]}{\partial s^{l}[i]}
\frac{\partial s^{l}[i]}{\partial u^{l}[i]}.
\end{equation}

\textbf{为什么有两项}：$u^l[i+1]$ 依赖 $u^l[i]$ 有两条路径：
\begin{itemize}[leftmargin=1.5em,itemsep=0pt]
\item \textbf{第一项} $\partial u^l[i+1]/\partial u^l[i]$：通过膜电位自身的 leak 传播。$u^l[i+1]=\lambda u^l[i]+\cdots$，所以 $\partial u^l[i+1]/\partial u^l[i]=\lambda I$；
\item \textbf{第二项} $\partial u^l[i+1]/\partial s^l[i]\cdot\partial s^l[i]/\partial u^l[i]$：通过 $u^l[i]$ 决定是否发放 $s^l[i]$，再通过 reset 项影响 $u^l[i+1]$。$\partial u^l[i+1]/\partial s^l[i]=-\lambda V_{th}I$（reset 项系数），$\partial s^l[i]/\partial u^l[i]=D_{\Psi,i}$（surrogate），所以该项 $=-\lambda V_{th}D_{\Psi,i}$。
\end{itemize}

\textbf{参数解释}：
\begin{itemize}[leftmargin=1.5em,itemsep=0pt]
\item $\epsilon^{l}[i]\in\mathbb R^{n^{(l)}\times n^{(l)}}$：第 $l$ 层从时间步 $i$ 到 $i+1$ 的 temporal Jacobian；
\item $\mathcal{E}_{t,\tau}$：从 $\tau$ 传播到 $t$ 的 ordered temporal product；
\item $\partial L/\partial s^{l}[t]\in\mathbb R^{n^{(l)}}$：第 $l$ 层脉冲 $s^l[t]$ 接收到的 same-time/spatial error signal（从输出层反向传播到本层）；
\item $\partial s^{l}[t]/\partial u^{l}[t]=D_{\Psi,t}$：surrogate derivative。
\end{itemize}

\textbf{Step 3：代入 LIF。}

把 $\partial u^l[i+1]/\partial u^l[i]=\lambda I$ 与 $\partial u^l[i+1]/\partial s^l[i]=-\lambda V_{th}I$ 代入 $\epsilon^l[i]$：
\begin{equation}
\epsilon^{l}[i]=\lambda I-\lambda V_{th}D_{\Psi,i}.
\end{equation}

\textbf{解释}：$\lambda I$ 来自 leak，$-\lambda V_{th}D_{\Psi,i}$ 来自 reset。$D_{\Psi,i}$ 是对角矩阵，其 $j$ 号对角元 $\Psi'(u_j-V_{th})$ 越接近 1，说明神经元 $j$ 处于阈值附近、是否发放对 $u$ 非常敏感，reset 项的贡献就越大。

\textbf{内存来源}：BPTT 在反向时需要 $\epsilon^l[i]$，而 $\epsilon^l[i]$ 依赖 $u^l[i]$ 与 $s^l[i]$。这意味着反向过程需要访问所有时间步的 forward states。典型实现要么把它们全部保存（内存 $O(TLn)$），要么在前向过程中 checkpoint 一些、反向时重算，总之无法做到 $O(1)$ 与 $T$ 无关的内存。

\subsection{STDP：局部突触可塑性基础}

\textbf{经典 STDP}。令 $\Delta t=t_{\mathrm{post}}-t_{\mathrm{pre}}$：
\begin{equation}
\Delta w=
\begin{cases}
A_+\,e^{-\Delta t/\tau_+}, & \Delta t>0\ \text{(potentiation)},\\[2mm]
-A_-\,e^{\Delta t/\tau_-}, & \Delta t<0\ \text{(depression)}.
\end{cases}
\end{equation}

\textbf{这个公式在说什么}：如果突触前神经元先发放、突触后神经元后发放（$\Delta t>0$），权重增强（potentiation）；反之如果突触后先发放、突触前后发放（$\Delta t<0$），权重减弱（depression）。这正是 Hebbian "fire together, wire together" 思想的时序版本。

\textbf{参数解释}：
\begin{itemize}[leftmargin=1.5em,itemsep=0pt]
\item $A_+,A_->0$：因果/非因果分支的振幅；
\item $\tau_+,\tau_->0$：两个分支的时间常数；$\Delta t$ 越大（越"不相关"），效应越弱；
\item 这个规则是严格\textbf{online}的：只需要当前和过去的信息，不需要未来。
\end{itemize}

\textbf{Generalized trace-based STDP form}（S-TLLR 所采用的形式）：
\begin{equation}
\Delta w_{ij}[t]
=
\alpha_{\mathrm{pre}}\,y_i[t]\sum_{t'=0}^{t}\lambda_{\mathrm{pre}}^{t-t'}x_j[t']
+
\alpha_{\mathrm{post}}\,x_j[t]\sum_{t'=0}^{t-1}\lambda_{\mathrm{post}}^{t-t'}y_i[t'].
\end{equation}

\textbf{这个公式在说什么}：把 STDP 从单次 pre/post 事件对推广到脉冲序列。第一项是\textbf{因果项}：当前时刻 post 发放 $y_i[t]=1$，与过去 pre 脉冲的加权和（权重随时间指数衰减）相关；第二项是\textbf{非因果项}：当前时刻 pre 发放 $x_j[t]=1$，与过去 post 脉冲的加权和（注意上界是 $t-1$，不包含当前 $y_i[t]$）相关。

\textbf{参数解释}：
\begin{itemize}[leftmargin=1.5em,itemsep=0pt]
\item $x_j[t]\in\{0,1\}$：pre 神经元 $j$ 在 $t$ 时刻的脉冲；
\item $y_i[t]\in\{0,1\}$：post 神经元 $i$ 在 $t$ 时刻的脉冲；
\item $\lambda_{\mathrm{pre}},\lambda_{\mathrm{post}}\in(0,1)$：pre/post 迹的指数衰减因子；与 LIF 的膜衰减 $\lambda$ 不同，它们是独立超参；
\item $\alpha_{\mathrm{pre}},\alpha_{\mathrm{post}}$：因果/非因果项的强度；$\alpha_{\mathrm{post}}=0$ 时只保留因果项。
\end{itemize}

\textbf{STDP 的局限}：它只使用 pre/post 时序，不包含任务误差信号。因此它本身不是监督学习中的完整 credit assignment 方案——它更新的是"最近同时活跃的突触"，而不是"让损失下降的突触"。

\textbf{从 STDP 到 three-factor rule}：引入任务信号 $\delta_i[t]$ 作为"第三个因子"来调制 eligibility：
\begin{equation}
\Delta w_{ij}[t]=\delta_i[t]\,e_{ij}[t].
\end{equation}

\textbf{解释}：$e_{ij}[t]$ 是由 STDP 时序决定的 eligibility（"该突触最近是否值得被修改"），$\delta_i[t]$ 是任务学习信号（"该修改是否有助于降低损失"）。两者相乘给出最终的权重更新。这个三因子结构是 S-TLLR 与 TESS 的基础。

\section{前向敏感度与局部资格迹：RTRL 与 e-prop}

\subsection{RTRL：forward-mode 参数敏感度}

\textbf{Step 1：一般 RNN 状态方程。}
\begin{equation}
h_t=F(h_{t-1},x_t,W).
\end{equation}
$h_t$ 是 $t$ 时刻的 hidden state（LIF 中 $h_t=u_t$），$F$ 是状态转移函数，$W$ 是所有可学习参数（flatten 后为 $r$ 维向量）。

\textbf{Step 2：参数敏感度。}
\begin{equation}
P_t:=\frac{\partial h_t}{\partial W}\in\mathbb R^{n\times r}.
\end{equation}

\textbf{这个量在说什么}：$P_t$ 的 $(i,k)$ 元素 $\partial h_{t,i}/\partial W_k$ 表示"第 $k$ 个参数变化一个单位，$t$ 时刻第 $i$ 个 hidden state 变化多少"。它包含了从初始时刻到 $t$ 时刻所有参数作用的累计影响。

\textbf{Step 3--4：链式法则与两个 Jacobian。}
\begin{equation}
\frac{\partial h_t}{\partial W}
=
\frac{\partial h_t}{\partial h_{t-1}}\cdot\frac{\partial h_{t-1}}{\partial W}
+
\left.\frac{\partial h_t}{\partial W}\right|_{h_{t-1}\text{ fixed}}.
\end{equation}

\textbf{解释}：$h_t$ 通过两条路径依赖 $W$——一是"经过 $h_{t-1}$"（即过去权重的累积影响被当前 dynamics 继续传播），二是"当前时刻 $W$ 直接进入 $F$"（$h_{t-1}$ 固定时 $W$ 对 $h_t$ 的直接作用）。定义 $A_t=\partial h_t/\partial h_{t-1}$、$B_t=\partial h_t/\partial W|_{h_{t-1}\text{ fixed}}$ 后：

\textbf{Step 5：RTRL 核心递推。}
\begin{equation}
P_t=A_t P_{t-1}+B_t.
\end{equation}

\textbf{这个公式在说什么}：这是一个\textbf{forward-in-time 递推}。$P_{t-1}$ 是过去所有参数作用的累计敏感度，乘以 $A_t$ 表示"过去的敏感度被当前 dynamics 继续传播"；再加上 $B_t$ 表示"当前时刻 $W$ 新注入的影响"。每前向一个时间步，就更新一次 $P_t$，不需要等整个序列完成。

\textbf{Step 6：当前 loss 梯度。} $\nabla_W\ell_t=\partial\ell_t/\partial h_t\cdot P_t$。

\textbf{参数解释}：
\begin{itemize}[leftmargin=1.5em,itemsep=0pt]
\item $P_t$：截至 $t$ 时刻累计的参数敏感度；
\item $A_t\in\mathbb R^{n\times n}$：hidden state 的自 Jacobian；
\item $B_t\in\mathbb R^{n\times r}$：参数的直接注入 Jacobian；
\item $\ell_t$：当前时刻的 loss；$\partial\ell_t/\partial h_t$ 是当前 loss 对 hidden state 的梯度。
\end{itemize}

\textbf{复杂度}：对 fully recurrent $n$-unit RNN、$r=O(n^2)$ 时，$P_t$ 需保存 $n\times r=O(n^3)$ 个元素，每步更新 $P_t=A_t P_{t-1}+B_t$ 需要 $O(n^4)$ 运算。多层结构的具体复杂度取决于 connectivity。

\textbf{与 BPTT 的关系}：BPTT 是 reverse-mode（从 $T$ 反向传播到 $1$），RTRL 是 forward-mode（从 $1$ 前向传播到 $T$）。对一个可微 computational graph 两者都是精确的，只是方向不同、内存分布不同。

\subsection{e-prop：从 RTRL 到局部资格迹 + 学习信号}

\textbf{Step 1：RTRL 的问题。} $P_t$ 的维度是 $n\times r$，对 SNN 来说太大（$n$ 是神经元数，$r$ 是权重总数）。e-prop 的核心想法是：\textbf{不直接保存 $P_t$，而是把它重新组织成"局部 eligibility $\times$ learning signal"}。

\textbf{Step 2：从经典 recurrent gradient factorization 出发。}
\begin{equation}
\frac{dE}{dW_{ji}}
=
\sum_{t'}
\frac{dE}{dh_j^{t'}}
\frac{\partial h_j^{t'}}{\partial W_{ji}}\Big|_{h_j^{t'-1}\text{ fixed}}.
\end{equation}

\textbf{这个公式在说什么}：损失对突触 $W_{ji}$ 的梯度，可以写成"对每个时间步 $t'$，hidden state $h_j^{t'}$ 对损失的贡献 $\times$ 该突触对 $h_j^{t'}$ 的直接影响"。

\textbf{Step 3：递归展开 $dE/dh_j^{t'}$。}

$dE/dh_j^{t'}$ 本身是 $h_j^{t'}$ 通过两条路径对 $E$ 的贡献：一是当前时刻 $h_j^{t'}$ 直接影响 $z_j^{t'}$（进而影响 $E$），二是 $h_j^{t'}$ 通过 $h_j^{t'+1}$ 影响未来。写成递归：
\begin{equation}
\frac{dE}{dh_j^{t'}}
=
L_j^{t'}\frac{\partial z_j^{t'}}{\partial h_j^{t'}}
+
\frac{dE}{dh_j^{t'+1}}\frac{\partial h_j^{t'+1}}{\partial h_j^{t'}},
\qquad
L_j^{t'}:=\frac{dE}{dz_j^{t'}}.
\end{equation}

\textbf{解释}：这是一个\textbf{反向递归}（从 $t'=T$ 往回算）。$L_j^{t'}=dE/dz_j^{t'}$ 是 ideal learning signal，包含所有从当前脉冲 $z_j^{t'}$ 出发、经其他神经元与后续时间步、最终影响 $E$ 的 downstream pathways。

把递归展开到 $T$、代回 Step 2、交换求和顺序（先对 $t$ 求和、再对 $t'$ 求和），可以重新组织成：
\begin{equation}
\frac{dE}{dW_{ji}}
=
\sum_t L_j^t\,e_{ji}^t,
\qquad
e_{ji}^t
:=
\left[\frac{dz_j^t}{dW_{ji}}\right]_{\text{local}}.
\end{equation}

\textbf{关于 $e_{ji}^t$ 的语义}：它是 chain-rule \textbf{rearrangement} 之后可由突触及 postsynaptic neuron 的局部 forward variables 计算的部分。其他 downstream 影响——$z_j^t$ 通过未来其他神经元、再回头影响 $E$ 的部分——\textbf{并未被丢弃}，而是全部进入了 $L_j^t$。因此这个分解是精确恒等式，不是"删掉某些路径后的近似"。

\textbf{Step 4：eligibility vector 递推。}

$e_{ji}^t$ 可以通过 eligibility vector $\xi_{ji}^t$ 前向递推得到：
\begin{equation}
e_{ji}^t=\frac{\partial z_j^t}{\partial h_j^t}\,\xi_{ji}^t,
\qquad
\xi_{ji}^{t+1}=\lambda\,\xi_{ji}^t+z_i^t.
\end{equation}

\textbf{解释}：$\xi_{ji}^t$ 是"滤波后的 presynaptic spike"，满足一个 leaky integrator 形式的递推。$\lambda$ 是 leak 系数，$z_i^t$ 是突触 $i\to j$ 的 pre 神经元 $i$ 在 $t$ 时刻的脉冲。$\partial z_j^t/\partial h_j^t$ 是 surrogate derivative（LIF 中等于 $\Psi'(u_j^t-V_{th})$）。

\textbf{Step 5：ideal → online → actual update。}

\textbf{(a) Ideal learning signal.} $L_j^t=dE/dz_j^t$ 是精确的、future-aware 的，但需要 BPTT 才能算出。

\textbf{(b) Online learning signal.} e-prop 用当前可得的信息近似它：

Regression：$\widehat L_j^t=\sum_k B_{jk}(y_k^t-y_k^{*,t})$；Classification：$\widehat L_j^t=\sum_k B_{jk}(\pi_k^t-\pi_k^{*,t})$。

\textbf{解释}：$y_k^t$ 是输出神经元 $k$ 的（滤波后）输出，$y_k^{*,t}$ 是目标；$\pi_k^t$ 是 softmax 概率；$B_{jk}$ 是 feedback 权重，表示"神经元 $j$ 到输出 $k$ 的误差路由权重"。三种选择：symmetric e-prop 用真实 readout 权重 $W_{kj}^{\mathrm{out}}$、random e-prop 用固定随机权重、adaptive e-prop 局部可塑性学习权重。

\textbf{(c) Actual synaptic update.} 由于 readout 是 leaky integrator，实际权重更新使用\textbf{滤波后的 eligibility} $\bar e_{ji}^t=\mathcal F_\kappa(e_{ji})^t$：
\begin{equation}
\Delta W_{ji}^{\mathrm{rec}}
=
-\eta\sum_t\Big[\sum_k B_{jk}(y_k^t-y_k^{*,t})\Big]\bar e_{ji}^t.
\end{equation}
注意被低通滤波的是 eligibility $\bar e_{ji}^t$，不是 error $y_k^t-y_k^{*,t}$。

\textbf{参数解释}：
\begin{itemize}[leftmargin=1.5em,itemsep=0pt]
\item $\lambda$（e-prop 原文中记作 $\alpha$）：postsynaptic state 的衰减因子；
\item $\xi_{ji}^t$：eligibility vector，保存突触 $i\to j$ 的局部历史；
\item $e_{ji}^t$：observable eligibility trace；
\item $\kappa$：readout dynamics 的低通衰减参数；
\item $B_{jk}$：output error 到 recurrent neuron $j$ 的 feedback weight；
\item $W_{kj}^{\mathrm{out}}$：神经元 $j$ 到 readout $k$ 的 forward weight；
\item $\eta$：学习率。
\end{itemize}

\textbf{与 RTRL 的关系}：RTRL 的 $P_t=A_tP_{t-1}+B_t$ 保存了完整网络的 Jacobian；e-prop 的 $\xi_{ji}^t=A_{j,t}^{\mathrm{loc}}\xi_{ji}^{t-1}+B_{ji,t}^{\mathrm{loc}}$ 只保留了单突触级的局部路径。二者递推结构相似，但对象完全不同——$P_t\in\mathbb R^{n\times r}$ 与 $\xi_{ji}^t\in\mathbb R$ 不是同维度、同语义的对象，不直接做数值距离比较。

\section{时间在线训练：OTTT 与 NDOT}

\subsection{从 BPTT 到 temporal dependency $\epsilon[i]$}

BPTT 的 temporal dependency：
\begin{equation}
\epsilon[i]
=
\frac{\partial u[i+1]}{\partial u[i]}
+
\frac{\partial u[i+1]}{\partial s[i]}\frac{\partial s[i]}{\partial u[i]}
=
\lambda I-\lambda V_{th}D_{\Psi,i}.
\end{equation}

\textbf{OTTT 与 NDOT 的区别就在于如何近似 $\epsilon[i]$}。

\subsection{OTTT：$\epsilon[i]\approx\lambda I$}

\textbf{Step 1：丢弃 spike-mediated path。} 注意到 Heaviside 几乎处处导数为零，$\partial s[i]/\partial u[i]=D_{\Psi,i}$ 在大多数时间步上很小，因此：
\begin{equation}
\frac{\partial u[i+1]}{\partial s[i]}\frac{\partial s[i]}{\partial u[i]}\approx 0.
\end{equation}

\textbf{这个近似在说什么}：忽略"$u[i]$ 通过 spike $s[i]$ 的 reset 影响 $u[i+1]$"这条路径。这样做的理由是——Heaviside 的导数在真正发放时可能很大，但在非发放时几乎为零，平均来看 $D_{\Psi,i}$ 对时间连乘的贡献远小于 $\lambda I$。

\textbf{Step 2：$\epsilon[i]$ 退化。}
\begin{equation}
\epsilon[i]\approx\lambda I.
\end{equation}

\textbf{Step 3：temporal trace。}

$\mathcal{E}_{t,\tau}=\epsilon[t-1]\cdots\epsilon[\tau]\approx\lambda^{t-\tau}I$。定义预突触活动的指数衰减迹：
\begin{equation}
\hat a^l[t]:=\sum_{\tau\le t}\lambda^{t-\tau}s^l[\tau],
\qquad
\hat a^l[t+1]=\lambda\hat a^l[t]+s^l[t+1].
\end{equation}

\textbf{这个公式在说什么}：$\hat a^l[t]$ 是对第 $l$ 层神经元从时刻 $1$ 到 $t$ 的所有历史脉冲 $s^l[\tau]$ 做指数加权求和。越近的脉冲权重越大（$\lambda^{t-\tau}$ 大），越早的脉冲权重越小。它满足一个简单的\textbf{前向递推}：$\hat a^l[t+1]$ 等于 $\hat a^l[t]$ 衰减一点、再加上 $s^l[t+1]$。这只需要存储一个向量，不需要保存历史 $s^l$。

\textbf{Step 4：instantaneous loss。} OTTT 为实现 online gradient，把传统依赖整段序列的 firing-rate loss 替换为逐时间步损失：
\begin{equation}
L[t]=\frac{1}{T}\mathcal{L}(s^N[t],y).
\end{equation}

\textbf{解释}：$s^N[t]$ 是输出层在 $t$ 时刻的脉冲，$y$ 是标签，$\mathcal{L}$ 是任务损失（如 CE）。每个时间步的损失独立计算，不再依赖整个序列。原论文明确指出：即使使用 instantaneous loss，OTTT 的 gradient 与使用 instantaneous loss 的 BPTT 在多层/循环网络上仍然不同，因为 OTTT 不保留 future influence。

\textbf{Step 5：OTTT 瞬时梯度。} 按 OTTT 原文的层索引 $W^l:l\to l+1$（即本文的 $W^{l+1\leftarrow l}$）：
\begin{equation}
\nabla_{W^l}L[t]=g_{u^{l+1}}[t]\,\hat a^l[t]^\top,
\qquad
g_{u^{l+1}}[t]=\left(\frac{\partial L[t]}{\partial s^{l+1}[t]}\frac{\partial s^{l+1}[t]}{\partial u^{l+1}[t]}\right)^\top.
\end{equation}

\textbf{这个公式在说什么}：当前时间步 $t$ 的权重梯度可以分解为"空间梯度 $\times$ 时间迹"的外积。$g_{u^{l+1}}[t]$ 是当前时刻 loss 对 $u^{l+1}[t]$ 的梯度（只需要当前时刻的 autograd，不需要跨时间反传）；$\hat a^l[t]$ 是历史 presynaptic spikes 的加权迹。

\textbf{Step 6：两种 update mode。}
\begin{equation}
\begin{array}{cl}
\text{OTTT}_O: & \text{每个 }t\text{ 立即用 }\nabla_{W^l}L[t]\text{ 更新}\\[1mm]
\text{OTTT}_A: & \text{累计 }\sum_t\nabla_{W^l}L[t]\text{ 后统一更新}
\end{array}
\end{equation}

\textbf{解释}：$O$ 表示 online，每个时间步更新一次参数；$A$ 表示 accumulated，把整段序列的梯度累加后再更新。前者更"在线"，后者在梯度保真比较中更接近 BPTT 的评估点。

\textbf{参数解释}：$\lambda$ 是膜衰减，同时成为 temporal trace decay；$\hat a^l[t]$ 是历史 presynaptic spikes 的加权迹；$g_{u^{l+1}}[t]$ 是当前时刻 spatial gradient；$T$ 是序列长度；$V_{th}$ 决定发放阈值并通过 surrogate derivative 影响 spatial gradient，但不进入 trace 的递推。

\textbf{复杂度}：额外 memory $O(Ln)$（每层存一个 trace 向量），per-step computation $O(Ln^2)$。

\subsection{NDOT：$\epsilon[i]\approx e[i]$}

\textbf{Step 1：OTTT 的问题。} OTTT 假设 $\epsilon[i]=\lambda I$，完全忽略了 spike 路径。但在某些时间步上，reset 路径的贡献可能并不小。NDOT 要保留 $\epsilon[i]$ 的动态性。

\textbf{Step 2：从连续动力学定义 $e(t)$。}

NDOT 引入隐式神经元动力学 $u(t+1)=\mathrm{Im}(u(t))$，对其求时间导数：
\begin{equation}
\frac{du(t+1)}{dt}=\frac{\partial u(t+1)}{\partial u(t)}\cdot\frac{du(t)}{dt}.
\end{equation}

由此定义 continuous temporal dependency：
\begin{equation}
e(t)
\triangleq
\frac{\partial u(t+1)}{\partial u(t)}
=
\frac{u'(t+1)}{u'(t)}.
\end{equation}

\textbf{这个公式在说什么}：$e(t)$ 是"下一时刻膜电位对当前时刻膜电位的偏导"，等于两个时刻膜电位变化率的比值。

\textbf{Step 3：代入 LIF 连续方程。} LIF 的连续膜电位方程是 $u'(t)=(-u(t)+I(t))/\tau$，其中 $I(t)$ 是输入电流，$\tau$ 是膜时间常数。代入得：
\begin{equation}
e(t)=\frac{u'(t+1)}{u'(t)}=\frac{u(t+1)-I(t+1)}{u(t)-I(t)}.
\end{equation}

\textbf{Step 4：映射到离散 SNN。} 将上述连续形式在离散时间步上取值，并结合论文采用的离散 LIF 更新方程，得到：
\begin{equation}
e^l[t]=
\begin{cases}
u^l[1]-V_{th}s^l[1], & t=1,\\[2mm]
\dfrac{u^l[t]-V_{th}s^l[t]}{u^l[t-1]-V_{th}s^l[t-1]}, & t>1.
\end{cases}
\end{equation}

\textbf{这个公式在说什么}：$e^l[t]$ 是"reset 后的膜电位"（$u^l[t]-V_{th}s^l[t]$）在相邻两个时刻的\textbf{逐元素比值}。$t=1$ 时用特殊初始化（NDOT 原文 Eq.(8) 后的 footnote 明确指出），避免 $u^l[0]-V_{th}s^l[0]$ 的退化处理。

需要强调：$e^l[t]$ 不是由离散 BPTT Jacobian $\epsilon^l[t]$ 严格代数化简得到的量。NDOT 的核心假设是：$e[t]$ 与 $\epsilon[t]$ 刻画相同的 temporal-dependency characteristic，因此可以用 $e[t]$ 近似/替代 $\epsilon[t]$。

\textbf{层索引关系}：NDOT 原文针对 $W^l$（即 $l-1\to l$）的具体梯度公式中，$\epsilon^l[i]$ 被 \textbf{presynaptic layer 的} $e^{l-1}[i]$ 替代（不是 $e^l[i]$），然后用逐元素乘法构造 $\hat a^{l-1}[t]$。

\textbf{Step 5：temporal trace。}
\begin{equation}
\hat a^{l-1}[t]=e^{l-1}[t-1]\odot\hat a^{l-1}[t-1]+s^{l-1}[t],
\end{equation}
展开为历史求和：
\begin{equation}
\hat a^{l-1}[t]
=
s^{l-1}[t]
+
\sum_{k<t}
\mathcal{E}^{\mathrm{NDOT}}_{t,k}\odot s^{l-1}[k],
\qquad
\mathcal{E}^{\mathrm{NDOT}}_{t,k}:=e^{l-1}[t-1]\odot\cdots\odot e^{l-1}[k].
\end{equation}

\textbf{解释}：与 OTTT 的结构完全相同，只是把固定标量 $\lambda$ 换成动态逐元素比率 $e^{l-1}[t]$。$\odot$ 是逐元素乘法，所以每个神经元有自己独立的时间衰减系数。

\textbf{Step 6：instantaneous loss。} NDOT 使用 CE 与 MSE 的加权和：
\begin{equation}
L[t]=(1-\alpha)\ell_{\mathrm{CE}}(s^N[t],y)+\alpha\,\ell_{\mathrm{MSE}}(s^N[t],y).
\end{equation}
$\alpha\in[0,1]$ 是权重系数，实验里通常取较小的值（让 CE 主导）。

\textbf{Step 7：NDOT 瞬时权重梯度与两种 update mode。}
\begin{equation}
\nabla_{W^l}L[t]=g_{u^l}[t]\,\hat a^{l-1}[t]^\top.
\end{equation}
这是当前时间步的 instantaneous weight gradient。完整累计更新为：
\begin{equation}
\begin{array}{cl}
\text{NDOT}_O: & \text{每个 }t\text{ 立即更新}\\[1mm]
\text{NDOT}_A: & \text{累计 }\sum_t\nabla_{W^l}L[t]\text{ 后统一更新}
\end{array}
\end{equation}

\textbf{参数解释}：
\begin{itemize}[leftmargin=1.5em,itemsep=0pt]
\item $e^l[t]\in\mathbb R^{n^{(l)}}$：neuron-wise dynamic temporal coefficient，逐元素比值；
\item $\hat a^{l-1}[t]\in\mathbb R^{n^{(l-1)}}$：temporal trace；
\item $g_{u^l}[t]\in\mathbb R^{n^{(l)}}$：当前时刻 spatial gradient（含从输出层反向传播到第 $l$ 层的完整 same-time spatial chain）；
\item $\odot$：逐元素乘法；
\item $I(t)$：连续膜电位方程中的输入电流；$\tau$：膜时间常数；
\item $V_{th}$：发放阈值；$\alpha$：CE 与 MSE 的权重；$T$：序列长度。
\end{itemize}

\textbf{数值稳定性（NDOT 原文 Appendix A.2）}：当 $e^l[t]$ 的分母 $u^l[t-1]-V_{th}s^l[t-1]$ 等于零时，NDOT 使用退化情况下的稳定策略：先根据膜电位符号处理，再用 clamping 把结果限制在 $[-\lambda,\lambda]$。这个 clamp \textbf{只对分母为零的退化情况生效}，不是 $e^l[t]$ 在所有时间步的一般定义。

\textbf{OTTT 与 NDOT 的关系}：
\begin{equation}
\begin{array}{cl}
\text{OTTT:} & \epsilon^l[t]\approx \lambda I\ \text{(matrix coefficient)}\\[1mm]
\text{NDOT:} & \epsilon^l[t]\approx \operatorname{diag}(e^l[t])\ \text{(diagonal dynamic coefficient)}
\end{array}
\end{equation}

\textbf{复杂度}：额外 memory $O(Ln)$，per-step computation $O(Ln^2)$，另加计算 $e^l[t]$ 的逐元素除法开销。

\section{STDP-inspired 本地学习：S-TLLR}

\textbf{Step 1：three-factor rule。} $\Delta w_{ij}[t]=\delta_i[t]\,e_{ij}[t]$。

\textbf{Step 2：eligibility。} 把 generalized trace-based STDP form 中的 post 脉冲 $y_i[t]$ 替换为 secondary activation $\Psi_{\mathrm{STLLR}}(u_i[t])$（$\Psi_{\mathrm{STLLR}}\neq\Psi'_{\mathrm{SG}}$）：
\begin{equation}
e_{ij}[t]
=
\alpha_{\mathrm{pre}}\Psi_{\mathrm{STLLR}}(u_i[t])\sum_{t'=0}^{t}\lambda_{\mathrm{pre}}^{t-t'}x_j[t']
+
\alpha_{\mathrm{post}}x_j[t]\sum_{t'=0}^{t-1}\lambda_{\mathrm{post}}^{t-t'}\Psi_{\mathrm{STLLR}}(u_i[t']).
\end{equation}

\textbf{这个公式在说什么}：第一项是因果项——当前 post 状态（用 $\Psi_{\mathrm{STLLR}}(u_i[t])$ 而非离散脉冲 $y_i[t]$ 表示）与过去 pre 脉冲的加权相关；第二项是非因果项——当前 pre 脉冲与过去 post 状态的加权相关。使用连续的 $\Psi_{\mathrm{STLLR}}(u_i[t])$ 而非二值 $y_i[t]$，可以让 eligibility 更平滑，便于学习。

\textbf{重要澄清：non-causal ≠ 偷看未来。} S-TLLR 所谓 "non-causal" 是 STDP timing relation 意义上的 non-causal——当前 pre spike 与\textbf{过去} post activity 的相关性。它仍然只使用当前和过去的信息，不需要访问未来时间步，因此仍是 online / temporally local。

\textbf{Step 3：前向递推形式。} 定义两个迹变量：
\begin{equation}
\operatorname{tr}(x_j)[t]=\lambda_{\mathrm{pre}}\operatorname{tr}(x_j)[t-1]+x_j[t],
\qquad
\operatorname{tr}(\Psi_{\mathrm{STLLR}}(u_i))[t]=\lambda_{\mathrm{post}}\operatorname{tr}(\Psi_{\mathrm{STLLR}}(u_i))[t-1]+\Psi_{\mathrm{STLLR}}(u_i[t]),
\end{equation}
\begin{equation}
e_{ij}[t]=\alpha_{\mathrm{pre}}\Psi_{\mathrm{STLLR}}(u_i[t])\operatorname{tr}(x_j)[t]+\alpha_{\mathrm{post}}x_j[t]\big(\operatorname{tr}(\Psi_{\mathrm{STLLR}}(u_i))[t]-\Psi_{\mathrm{STLLR}}(u_i[t])\big).
\end{equation}
括号里减去 $\Psi_{\mathrm{STLLR}}(u_i[t])$ 是因为 non-causal 项的求和上界是 $t-1$，不包含当前时刻。

\textbf{Step 4：recurrent synapse extension。} 对 recurrent connection：
\begin{equation}
u_i[t]=\gamma(u_i[t-1]-v_{th}y_i[t-1])+W^{\mathrm{ff}}x[t]+W^{\mathrm{rec}}y[t-1],
\end{equation}
presynaptic activity 变成上一时刻的 $y_k[t-1]$，eligibility 相应调整为：
\begin{equation}
e_{ik}^{\mathrm{rec}}[t]=\alpha_{\mathrm{pre}}\Psi_{\mathrm{STLLR}}(u_i[t])\sum_{t'=1}^{t}\lambda_{\mathrm{pre}}^{t-t'}y_k[t'-1]+\alpha_{\mathrm{post}}y_k[t-1]\sum_{t'=0}^{t-1}\lambda_{\mathrm{post}}^{t-t'}\Psi_{\mathrm{STLLR}}(u_i[t']).
\end{equation}

\textbf{Step 5--6：learning signal 与更新。}
\begin{equation}
w_{ij}:=w_{ij}+\rho\sum_{t=T_l}^{T}\delta_i[t]\,e_{ij}[t].
\end{equation}

\textbf{参数解释}：
\begin{itemize}[leftmargin=1.5em,itemsep=0pt]
\item $\Psi_{\mathrm{STLLR}}$：secondary activation function，与 firing function 不同，不必等于 $\Psi'_{\mathrm{SG}}$；
\item $\alpha_{\mathrm{pre}},\alpha_{\mathrm{post}}$：因果/非因果项强度；$\alpha_{\mathrm{post}}=0$ 只保留因果；
\item $\lambda_{\mathrm{pre}},\lambda_{\mathrm{post}}\in(0,1)$：因果/非因果迹的衰减因子；
\item $x_j[t],y_i[t]$：pre/post 脉冲；
\item $T_l$：learning signal 起始时间步，$t<T_l$ 时 $\delta_i[t]=0$；典型值为最后几步；
\item $\rho$：学习率；
\item recurrent extension 中 $\gamma,v_{th},W^{\mathrm{ff}},W^{\mathrm{rec}},y_k$ 分别为膜衰减、发放阈值、feedforward weights、recurrent weights、recurrent presynaptic spike。
\end{itemize}

\textbf{复杂度}：额外 memory $O(Ln)$（瞬时 eligibility，不保存 synapse-wise 状态），per-step computation $O(Ln^2)$。

\textbf{与 BPTT 的区别}：S-TLLR 显式引入 STDP-inspired \textbf{non-causal timing component}；该 component \textbf{不是由 BPTT 的 causal temporal chain-rule gradient 推导得到的}。BPTT 及其近似只利用 causal relations，而 S-TLLR 显式增加 causal + non-causal spike-timing relations。

\textbf{与 e-prop 的区别}：S-TLLR 的 eligibility 不要求 $\lambda$ 与膜衰减一致，$\lambda_{\mathrm{pre}},\lambda_{\mathrm{post}}$ 是独立超参。

\section{时空全局部学习：TESS}

\textbf{TESS 与 S-TLLR 的关系}：
\begin{equation}
\begin{array}{cl}
\text{S-TLLR:} & \text{local temporal eligibility} + \text{instantaneous BP learning signal}\\
\text{TESS:} & \text{local temporal eligibility} + \text{explicitly layer-local generated learning signal}
\end{array}
\end{equation}

\textbf{解释}：S-TLLR 的 temporal 部分是局部的（eligibility 只依赖 pre/post），但 spatial 部分仍然要跨层 BP 生成 learning signal $\delta_i[t]$。TESS 的推进是把 learning signal 也本地化——每层通过固定 basis matrix $B^{(l)}$ 独立生成，无需跨层 BP。

\textbf{Step 1--2：因果迹与非因果迹。}
\begin{equation}
q^{(l)}[t]=\lambda_{\mathrm{pre}}q^{(l)}[t-1]+o^{(l-1)}[t],
\qquad
h^{(l)}[t]=\lambda_{\mathrm{post}}h^{(l)}[t-1]+\Psi_{\mathrm{STLLR}}(u^{(l)}[t-1]).
\end{equation}
$q^{(l)}$ 是 presynaptic trace（因果迹），$h^{(l)}$ 是 postsynaptic activation trace（非因果迹）。

\textbf{Step 3：本地学习信号（LSG）。}
\begin{equation}
m^{(l)}[t]=B^{(l)\top}\big(f(B^{(l)}o^{(l)}[t])-y^*\big).
\end{equation}

\textbf{这个公式在说什么}：把本层输出 $o^{(l)}[t]$ 用固定二值矩阵 $B^{(l)}$ 投影到任务空间（$C$ 维），应用 softmax/identity 得到预测，与标签 $y^*$ 比较得到误差，再投影回来。这个 $m^{(l)}[t]$ 是\textbf{本层自己产生的} learning signal，不需要从其他层反传。

\textbf{Step 4：逐时刻权重增量。}
\begin{equation}
\Delta W^{(l)}[t]
=
\big(m^{(l)}[t]\odot\alpha_{\mathrm{pre}}\Psi_{\mathrm{STLLR}}(u^{(l)}[t])\big)\otimes q^{(l)}[t]
+
\big(m^{(l)}[t]\odot\alpha_{\mathrm{post}}h^{(l)}[t]\big)\otimes o^{(l-1)}[t].
\end{equation}
第一项是因果贡献（pre trace $q$ 与本层 activation 的外积），第二项是非因果贡献（本层 post trace $h$ 与 pre 输出的外积）。

\textbf{Step 5：最终权重更新与初始化。}
\begin{equation}
W^{(l)}\leftarrow W^{(l)}+\eta\sum_{t=t_l}^{T}\Delta W^{(l)}[t],
\qquad
q^{(l)}[0]=0,\quad h^{(l)}[0]=0.
\end{equation}
且 $m^{(l)}[t]$ 与 $\Delta W^{(l)}[t]$ 只在 $t\ge t_l$ 时计算。

\textbf{参数解释}：
\begin{itemize}[leftmargin=1.5em,itemsep=0pt]
\item $B^{(l)}\in\mathbb R^{C\times n^{(l)}}$：固定二值投影矩阵，列为方波函数，各列准正交；不参与训练；$C$ 是类别数，$n^{(l)}$ 是第 $l$ 层神经元数；
\item $o^{(l)}[t]\in\mathbb R^{n^{(l)}}$：第 $l$ 层 $t$ 时刻的输出脉冲；
\item $q^{(l)}[t]\in\mathbb R^{n^{(l-1)}}$：presynaptic trace（因果迹），维度由输入神经元数决定；
\item $h^{(l)}[t]\in\mathbb R^{n^{(l)}}$：postsynaptic activation trace（非因果迹），维度由本层神经元数决定；
\item $m^{(l)}[t]\in\mathbb R^{n^{(l)}}$：本层本地生成的 learning signal；
\item $f$：softmax（分类）或 identity（回归）；
\item $\Psi_{\mathrm{STLLR}}$：secondary activation，与 $\Psi'_{\mathrm{SG}}$ 不同；
\item $\lambda_{\mathrm{pre}},\lambda_{\mathrm{post}}$：因果/非因果迹衰减；
\item $\alpha_{\mathrm{pre}},\alpha_{\mathrm{post}}$：因果/非因果强度；
\item $t_l$：开始生成 learning signal 的时刻；$\eta$：学习率。
\end{itemize}

\textbf{形状检查}：$\Psi_{\mathrm{STLLR}}(u^{(l)}[t])\in\mathbb R^{n^{(l)}}$，$q^{(l)}[t]\in\mathbb R^{n^{(l-1)}}$，$\Psi_{\mathrm{STLLR}}(u^{(l)}[t])\otimes q^{(l)}[t]\in\mathbb R^{n^{(l)}\times n^{(l-1)}}$，与 $W^{(l)}$ 形状一致。

\textbf{复杂度}：额外 memory $O(Ln)$，learning-signal generation computation $O(LCn)$（$C$ 为类别数，$C\ll n$）。

\textbf{与 S-TLLR 的区别}：S-TLLR 的 $\delta_i[t]$ 由跨层 BP 生成（也可用 DFA）；TESS 的 $m^{(l)}[t]$ 完全在层内生成。

\section{算法关系与核心研究问题}

\subsection{三条研究路径}

\begin{equation}
\begin{array}{ccccccc}
& \textbf{BPTT} & & & & &\\
& \epsilon[t] & & & & &\\
\swarrow & & \searrow & & & &\\
\textbf{OTTT} & & \textbf{NDOT} & & & &\\
\lambda & & e[t] & & & &\\[3mm]
\textbf{RTRL} & \longrightarrow & \textbf{e-prop} & & & &\\
P_t & & L_j^t e_{ji}^t & & & &\\[3mm]
\textbf{STDP} & \longrightarrow & \textbf{S-TLLR} & \longrightarrow & \textbf{TESS} & &\\
\text{pre/post} & & \delta e & & m^{(l)}e & &
\end{array}
\end{equation}

\textbf{关于 RTRL $\to$ e-prop 箭头}：这是本文的研究比较路径，不意味着 e-prop 是直接由 RTRL 的 $P_t$ 局部截断所得。

\subsection{路径一：BPTT $\to$ OTTT $\to$ NDOT（时间近似）}

核心问题：$\epsilon^l[t]\to\lambda I$ 或 $\operatorname{diag}(e^l[t])$ 改变了多少 temporal credit？

\textbf{实验分两类}：

\begin{itemize}[leftmargin=1.5em,itemsep=0pt]
\item \textbf{A. Paper-faithful reproduction}：各方法按原论文 loss / update mode 跑。
\item \textbf{B. Controlled mechanism comparison}：统一 controlled instantaneous objective
\begin{equation}
\ell_{\mathrm{ctrl}}[t]=m[t]\,\ell_{\mathrm{CE}}\big(r[t],y[t]\big),
\qquad
L_{\mathrm{ctrl}}=\sum_{t=1}^{T}\ell_{\mathrm{ctrl}}[t].
\end{equation}
BPTT、OTTT、NDOT 三个 gradient estimator 全部基于同一个 $L_{\mathrm{ctrl}}$。
\end{itemize}

\textbf{受控比较的唯一变量} = temporal dependency treatment：
\begin{equation}
\epsilon^l[t]\quad\text{vs}\quad\lambda I\quad\text{vs}\quad\operatorname{diag}(e^l[t]).
\end{equation}

\textbf{Representation fidelity（同层比较）}：
\begin{equation}
D_{\epsilon,\lambda}^{l}[t]
=
\frac{\|\epsilon^l[t]-\lambda I\|_F}{\|\epsilon^l[t]\|_F+\varepsilon_{\mathrm{num}}},
\qquad
D_{\epsilon,e}^{l}[t]
=
\frac{\|\epsilon^l[t]-\operatorname{diag}(e^l[t])\|_F}{\|\epsilon^l[t]\|_F+\varepsilon_{\mathrm{num}}},
\qquad t\ge 2.
\end{equation}

\textbf{Weight-gradient 层面}：
\begin{equation}
D_{\mathrm{OTTT}}=\frac{\|g_{\mathrm{OTTT}}-g_{\mathrm{BPTT}}\|}{\|g_{\mathrm{BPTT}}\|+\varepsilon_{\mathrm{num}}},
\qquad
D_{\mathrm{NDOT}}=\frac{\|g_{\mathrm{NDOT}}-g_{\mathrm{BPTT}}\|}{\|g_{\mathrm{BPTT}}\|+\varepsilon_{\mathrm{num}}},
\end{equation}
\begin{equation}
\cos(g_{\mathrm{OTTT}},g_{\mathrm{BPTT}}),\qquad\cos(g_{\mathrm{NDOT}},g_{\mathrm{BPTT}}).
\end{equation}

\textbf{One-step loss change}：对每个 gradient estimator $\hat g$，在固定 $W$ 处测量一步更新后的损失变化：
\begin{equation}
\Delta L(\hat g)
:=
L_{\mathrm{ctrl}}(W-\eta\hat g)-L_{\mathrm{ctrl}}(W).
\end{equation}
对 $g_{\mathrm{BPTT}},g_{\mathrm{OTTT}},g_{\mathrm{NDOT}}$ 分别计算并比较。
\begin{itemize}[leftmargin=1.5em,itemsep=0pt]
\item 它回答的问题比 $\cos$ 更直接：方向大致对，但是否真的导致 loss 下降？
\item $\Delta L<0$ 表示该 estimator 在当前点仍然是一个 descent direction；
\item $\Delta L\ge 0$ 表示该 estimator 可能不是有效的下降方向，即使 $\cos(g,g_{\mathrm{BPTT}})>0$；
\item 这一步是 gradient similarity 与 final accuracy 之间的桥梁。
\end{itemize}

\textbf{多步优化与任务性能}：由 one-step 结果再推导 multi-step training trajectory 与最终 task accuracy 是否被 same interaction structure 解释。

\subsection{路径二：RTRL $\to$ e-prop（精确到局部）}

核心问题：RTRL 如何在 forward mode 保存完整网络 sensitivity？e-prop 如何通过 local eligibility + learning signal 重新组织相同的 gradient credit？

研究思路：tiny recurrent network 上实现并验证 RTRL 递推；独立实现 e-prop eligibility；比较 memory、computation、locality；研究 ideal $L_j^t$ 替换为 online $\widehat L_j^t$ 后 gradient error 多大。

\subsection{路径三：STDP $\to$ S-TLLR $\to$ TESS（生物局部化）}

\textbf{单因素 ablation}：

\textbf{Ablation 1：固定 temporal trace，替换 learning signal}。$\text{fixed trace}\times\{\delta_{\mathrm{BP}}[t],m_{\mathrm{TESS}}^{(l)}[t]\}$。

\textbf{Ablation 2：固定 learning signal，替换 trace}。$\text{fixed signal}\times\{\text{causal-only},\text{causal+non-causal}\}$。

\section{符号约定与复杂度口径}

\subsection{符号统一}

以下三个 $e$ 是不同对象：
\begin{equation}
\begin{array}{cl}
e^{l}_{\mathrm{NDOT}}[t]\in\mathbb R^{n^{(l)}} & \text{神经元动力学 temporal coefficient}\\[1mm]
e^{t}_{ji,\mathrm{eprop}}\in\mathbb R & \text{突触级 eligibility trace}\\[1mm]
e_{ij,\mathrm{STLLR}}[t]\in\mathbb R & \text{STDP-inspired eligibility}
\end{array}
\end{equation}

e-prop 原论文将其 eligibility vector 记为 $\epsilon_{ji}^t$；本文重命名为 $\xi_{ji}^t$，避免与 $\epsilon^l[t]$ 混淆。

\textbf{两种 $\Psi$ 也要区分}：
\begin{equation}
\Psi'_{\mathrm{SG}}(u)\neq\Psi_{\mathrm{STLLR}}(u).
\end{equation}
BPTT 中的 $\Psi'_{\mathrm{SG}}$ 是 spike surrogate derivative；S-TLLR/TESS 中的 $\Psi_{\mathrm{STLLR}}$ 是 secondary activation function。

\subsection{层索引与论文原下标映射}

OTTT 原论文 $W^l$ 表示 $l\to l+1$；NDOT 原论文 $W^l$ 表示 $l-1\to l$。本文基础部分统一使用 $W^{l\leftarrow l-1}$ 表示从第 $l-1$ 层到第 $l$ 层的权重。

\subsection{复杂度统一口径}

复杂度仅在相同网络假设下直接比较。

\begin{center}
\begin{tabular}{lll}
\toprule
算法 & 额外 memory & per-step computation \\
\midrule
BPTT        & $O(TLn)$              & $O(Ln^2)$ \\
RTRL        & $O(n^3)$              & $O(n^4)$ \\
e-prop      & $O(Ln^2)$             & $O(Ln^2)$ \\
OTTT        & $O(Ln)$               & $O(Ln^2)$ \\
NDOT        & $O(Ln)$               & $O(Ln^2)$ + ratio computation \\
S-TLLR      & $O(Ln)$               & $O(Ln^2)$ \\
TESS        & $O(Ln)$               & $O(LCn)$ \\
\bottomrule
\end{tabular}
\end{center}

\section{下一步实验与实现方案}

\subsection{实现顺序（Level 0 / 1 / 2）}

\textbf{Level 0：必须自己实现}

\begin{itemize}[leftmargin=1.5em,itemsep=0pt]
\item \textbf{BPTT-SG reference estimator}：作为 oracle，提供 $g_{\mathrm{BPTT}}$；
\item \textbf{Temporal-local primitive}：OTTT 风格的固定衰减 trace（$\hat a^l[t]=\lambda\hat a^l[t-1]+s^l[t]$）；
\item \textbf{Spatial-local primitive}：Mostafa-style fixed random local classifier，每层独立 local objective。
\end{itemize}

这是实验真正需要控制的三个核心组件。

\textbf{Level 1：必须验证}

\begin{itemize}[leftmargin=1.5em,itemsep=0pt]
\item \textbf{e-prop}：作为 eligibility-trace / learning-signal 分解的经典理论基线；实现 ideal factorization 与 online $\widehat L_j^t$ 两种。
\end{itemize}

\textbf{Level 2：external validation（优先使用作者代码）}

\begin{itemize}[leftmargin=1.5em,itemsep=0pt]
\item \textbf{OTTT} 与 \textbf{NDOT}：官方实现已公开，优先使用原代码；
\item \textbf{S-TLLR} 与 \textbf{TESS}：TESS 官方仓库已区分 S-TLLR baseline 与 TESS full；
\item \textbf{Traces Propagation} 等后续工作：作为额外 external check，不要求从零重写。
\end{itemize}

不需要在第一阶段把五个算法全部从零重写。

\subsection{基础验证：按三条论文路径组织}

\textbf{路径一}：验证 BPTT 手写梯度与 autograd 一致；验证 OTTT 的 trace 递推与闭式求和一致；验证 NDOT 的 trace 递推与历史展开式一致。

\textbf{路径二}：验证 RTRL 递推与 autograd full Jacobian 一致；验证 ideal e-prop identity；验证 eligibility vector 局部递推；再单独研究 online learning signal。

\textbf{路径三}：验证 S-TLLR 的显式历史求和与 recursive trace 实现一致；验证 TESS 的 $m^{(l)}[t]$ 与最终 $W^{(l)}\leftarrow W^{(l)}+\eta\sum_{t=t_l}^{T}\Delta W^{(l)}[t]$。

\subsection{Synthetic task 定义与 readout}

\textbf{Readout 定义（所有实验统一）}：
\begin{equation}
r[t]=W_{\mathrm{out}}s^N[t]+b_{\mathrm{out}},
\qquad
W_{\mathrm{out}}\in\mathbb R^{C\times n^N},\ b_{\mathrm{out}}\in\mathbb R^C.
\end{equation}

\textbf{Readout 参数在 gradient fidelity diagnostic 阶段固定，不参与训练。}这确保被比较的 gradient estimator 是关于同一组被训练参数（hidden / feedforward weights）的。

\textbf{Delayed Copy}：
\begin{itemize}[leftmargin=1.5em,itemsep=0pt]
\item 输入：$x[1]\in\{0,1\}$，$x[2],\dots,x[k]=0$；
\item 目标：$y[k+1]=x[1]$；
\item mask：$m[k+1]=1$，其余 $m[t]=0$。
\end{itemize}

\textbf{Delayed XOR}：
\begin{itemize}[leftmargin=1.5em,itemsep=0pt]
\item 输入：$x[1]=b_1$，$x[2]=b_2$；之后 $x[t]=0$；
\item 目标：$y[k+2]=b_1\oplus b_2$；
\item mask：$m[k+2]=1$，其余 $m[t]=0$。
\end{itemize}

\textbf{关于 Delayed XOR 中 $k$ 的定义}：$k$ 是 $b_2$（最后一个 informative input）与 target 之间的 delay。从 $x[1]$ 到 target 实际是 $k+1$。

\subsection{第一组实验}

\textbf{路径一}：单层。参数网格：
\begin{equation}
T\in\{4,8,16,32,64\},
\qquad
k\in\{1,2,4,8,16,32\}\cap[1,T-1].
\end{equation}

\textbf{建议实验设计}：
\begin{itemize}[leftmargin=1.5em,itemsep=0pt]
\item 固定 $T=64$ 扫描 $k\in\{1,2,4,8,16,32\}$；
\item 固定 $k=8$ 扫描 $T\in\{16,32,64\}$。
\end{itemize}

\textbf{Gradient fidelity 测量}：默认使用 A-mode（frozen-$W$ diagnostic）。BPTT oracle 使用同一 controlled loss $L_{\mathrm{ctrl}}$。

绘制 $D_{\mathrm{OTTT}}(k)$、$D_{\mathrm{NDOT}}(k)$、$\cos$ 曲线、$\Delta L(k)$ 曲线，以及同层 representation fidelity。

\textbf{路径二}：tiny recurrent network 上比较 RTRL 与 e-prop 的 memory、computation、locality 与 gradient error。

\textbf{路径三}：单因素 ablation。

\subsection{后续扩展原则与 decision criterion}

后续研究方向不预先固定。首先完成上述三条研究路径，依据实验发现的 failure mode 决定后续方向。

\textbf{判断后续是否进入机制设计，使用 effect-size / statistical criterion，不用 accuracy 百分比阈值}：

\begin{itemize}[leftmargin=1.5em,itemsep=0pt]
\item \textbf{Gradient-level}: $\Delta\cos(\hat g,g_{\mathrm{BPTT}})$ 是否稳定且跨任务复现；
\item \textbf{Optimization-level}: one-step loss change $\Delta L(\hat g)$ 是否具有稳定的符号与量级；
\item \textbf{Task-level}: task accuracy 的变化是否与 gradient-level / optimization-level 的效应一致。
\end{itemize}

任何 single-number 的 accuracy 阈值（例如 $1\%$）都不构成科学判据。

\textbf{新机制的引入原则}：只有在存在稳定、可重复、真正影响优化或资源效率的 failure mode 后，才进一步设计新算法。

\bibliographystyle{IEEEtran}
\bibliography{SNN_OnlineLearningSummary}
\end{document}